\scanpage{49}
En fait, nous allons travailler avec des éléments de $\Pi_1(\widetilde\Sigma)$ qui proviennent, soit des p\textsuperscript{ts} de $\widetilde\Sigma$, donc des p\textsuperscript{ts} de $\Sigma^*$ ou d'un des $U_i$, soit des \uline{parties} de $\widetilde\Sigma$ qui sont 1-connexes pour la top. induite par $\widetilde\Sigma$ -- en pratique peut-être des parties \inserted{1-conn.} de $\Sigma^*$, et en les faisant ``déborder'' sur $\widetilde\Sigma$ de façon naturelle.

Le groupe $\Gamma$ opère sur la situation \inserted{$C^\infty$} $(\Sigma,J)$, donc par transport de structure sur la situation
\[
\widetilde\Sigma\longrightarrow\Sigma,\quad J\subset\Sigma
\]
donc sur \deleted{le} (14) (15). Un aspect essentiel du travail à faire est de garder \deleted{cette} constamment présente à l'esprit cette opération, et d'introduire des notations et des constructions qui soient invariantes, dans le sens requis par la situation, par les opérations en question.
\margin{48}
